\documentclass[11pt,a4paper]{article}
\usepackage[T1]{fontenc}
\usepackage{lmodern,amsmath,amssymb,amsthm,mathtools,booktabs,microtype}
\usepackage[margin=27mm]{geometry}
\usepackage[colorlinks=true,linkcolor=blue,urlcolor=blue]{hyperref}
\setlength{\emergencystretch}{2em}
\newtheorem{theorem}{Theorem}[section]
\newtheorem{lemma}[theorem]{Lemma}
\newtheorem{corollary}[theorem]{Corollary}
\newcommand{\Q}{\mathbb Q}\newcommand{\Z}{\mathbb Z}\newcommand{\F}{\mathbb F}
\title{Mixed torsion signs for doubled elliptic reference points}
\author{}\date{16 November 2026}
\begin{document}\maketitle
\begin{abstract}
An explicit progression of doubled points on the project elliptic curve
reduces to the identity at one prime above a split rational prime, while
its translate reduces to the distinguished two-torsion point at the
conjugate prime. The two resulting formal-parameter signs are opposite.
The construction stays inside the complete reference chart and gives
infinitely many characteristic-zero points. Exact reduction orders and
an independent finite-field certificate prove every assertion. The
result excludes automatic agreement of these signs, including for
doubled reference points; it makes no assertion about solutions of a
complete detector system.
\end{abstract}

\section{An explicit progression}
Put
\[
 K=\Q(\rho),\qquad \rho^2=2,\qquad D=11-6\rho,
 \qquad E:y^2=x^3-Dx,
\]
\[
 R_0=(1+\rho,2\rho),\qquad
 P=2R_0=\left(\frac{57}{8}-\frac72\rho,
                  -\frac{165}{8}+\frac{417}{32}\rho\right),
 \qquad T_0=(0,0).
\]
For a rational prime $p$ and a root $r$ of $r^2=2$ in $\F_p$,
reduction at $(p,\rho-r)$ is denoted by $\operatorname{red}_{p,r}$.
At a prime of good reduction this denotes the group homomorphism
on all of $E(K)$, including points reducing to the identity $O$.

\begin{theorem}\label{thm:progression}
Let
\[
 p=76537,\qquad r_1=26935,\qquad r_2=49602,
\]
and for every integer $j\ge0$ set
\[
 h_j=28101000+30510750j,\qquad H_j=[h_j]P.
\]
Then $p$ is prime, $r_1,r_2$ are the two roots of $X^2-2$ in
$\F_p$, and $E$ has good reduction at both corresponding primes.
The points $H_j$ are distinct, belong to $2E(K)$, and satisfy
\[
 \operatorname{red}_{p,r_1}(H_j)=O,
 \qquad
 \operatorname{red}_{p,r_2}(H_j-P)=T_0.
\]
For every $j$, the points $H_j,H_j-P$ are finite and have nonzero
abscissas. Thus $H_j$ belongs to
\[
                     E\setminus\{O,T_0,P,P+T_0\},
\]
the full domain of the three-coordinate reference chart.
\end{theorem}

\section{Finite reduction certificates}
The following data determine the two reductions and their exact point
orders. All entries in this section are ordinary residues modulo the
indicated prime.
\[
\begin{array}{c|r|r|r}
 \rho&D&P&\operatorname{ord}(P)\\\hline
 26935&68012&(10973,56787)&6375\\
 49602&8547 &(46444,38843)&4786
\end{array}
\]
The integer $76537$ has no divisor from $2$ through
$\lfloor\sqrt{76537}\rfloor=276$. Both displayed values of $D$
are nonzero, so the discriminant $64D^3$ is nonzero modulo $p$.
The two listed roots of $2$ have sum $76537$, as required for the
conjugate primes.

The exact orders have short certificates. In the first reduction,
$[6375]P=O$ and
\[
\begin{array}{c|r}
 \text{multiplier}&\text{point}\\\hline
 6375/3 &(7526,20385)\\
 6375/5 &(5801,44942)\\
 6375/17&(52213,15802)
\end{array}
\]
are all nonzero. The prime divisors of $6375$ are precisely $3,5,17$.
In the second reduction, $[4786]P=O$ and
\[
 [4786/2]P=(0,0),\qquad
 [4786/2393]P=(46317,70872),
\]
where $4786=2\cdot2393$ and $2393$ is prime. These identities prove
the orders without an assumption about the full finite groups.
They also give
\begin{equation}\label{eq:half}
                   [2393]\operatorname{red}_{p,r_2}(P)=T_0.
\end{equation}

\begin{lemma}\label{lem:nontorsion}
The point $P$ has infinite order.
\end{lemma}
\begin{proof}
At the good prime $(17,\rho-6)$, the curve has $D=9$ and
$P=(1,14)$; this reduction has exact order $5$.
Suppose $P$ had finite order. Its reduction of order $6375$ above
shows that its order has a nontrivial $3$-primary part.
A nonzero point in that $3$-primary subgroup would reduce to $O$
at $(17,\rho-6)$, because its reduced order divides both a power
of $3$ and $5$.

The kernel of good reduction at $17$ has no nonzero $3$-power
torsion. Indeed its formal parameter $t$ has positive valuation and
satisfies $[3](t)=3t+O(t^2)$ with integral coefficients. Since $3$
is a $17$-adic unit, $[3](t)$ has the same finite valuation as $t$.
Iteration cannot give zero. This contradiction proves the lemma.
\end{proof}

\begin{proof}[Proof of Theorem~\ref{thm:progression}]
The finite arithmetic just given verifies primality, splitting and good
reduction. The exact orders are coprime and their product is
$30510750$. Moreover
\[
 28101000=6375\cdot4408,\qquad
 28101000-1=4786\cdot5871+2393.
\]
These congruences hold for every $h_j$. Reduction commutes with
integer multiplication, so the first gives the identity at the first
prime, while the second and~\eqref{eq:half} give $T_0$ at the second.

The equality $P=2R_0$ gives $H_j=[2h_j]R_0\in2E(K)$.
Lemma~\ref{lem:nontorsion} makes the $H_j$ distinct.
Both $h_j$ and $h_j-1$ are nonzero. Thus $H_j$ and $H_j-P$ are
non-torsion, so neither is $O$ or $T_0$. On this Weierstrass equation,
the only point of abscissa zero is $T_0$. Both points are therefore
finite with nonzero abscissas, proving the final assertion.
\end{proof}

\section{The formal-parameter signs}
Use the local parameter
\[
                         \lambda(Q)=-x(Q)/y(Q)
\]
at $O$. Translation by $T_0$ is given, whenever $x\ne0$, by
\begin{equation}\label{eq:translation}
 (x,y)\longmapsto
       \left(-\frac D x,\frac{Dy}{x^2}\right).
\end{equation}
The chord formula proves this identity directly. In particular,
\begin{equation}\label{eq:sign}
 \frac{y(Q+T_0)}{x(Q+T_0)}=-\frac{y(Q)}{x(Q)},
 \qquad
                         \lambda(Q+T_0)=-\lambda(Q).
\end{equation}
These are rational-function identities, including their regular
extensions at the relevant formal neighborhoods.

\begin{corollary}\label{cor:signs}
For every $H_j$ of Theorem~\ref{thm:progression}, the primary reference
$H_j$ is a nonidentity point in the formal group at $O$ for the first
completion. For the second completion, the nonidentity formal point
is $(H_j-P)+T_0$, and
\[
             \lambda(H_j-P)=-\lambda((H_j-P)+T_0).
\]
Hence identifying the two reference neighborhoods with the formal group
at $O$ produces opposite translation signs. The signs need not agree
for a single global doubled reference point.
\end{corollary}
\begin{proof}
The reduction assertions put the indicated points in the formal groups.
They are nonidentity because their characteristic-zero points are
non-torsion. The first identification uses no torsion translation.
The second uses $T_0$, and its minus sign is exactly~\eqref{eq:sign}.
\end{proof}

The corollary concerns the reference points and their local parameters.
It does not supply witnesses for any contact, ratio, trace or pole
system. In particular it is not a counterexample to an established
defining polynomial. It rules out deducing automatic agreement of these
two signs from the reference-chart equations and membership in $2E(K)$
alone. A construction may impose further conditions that exclude or
control the mixed branches; those conditions require their own proof.

\section{Reproducible exact verification}
The directory
\begin{center}\small
\path{negative_results/computation/mixed_sign_v42}
\end{center}
contains the complete finite certificate. The producer uses affine
addition and right-to-left binary multiplication. The checker imports
neither the producer nor a symbolic algebra package: it reconstructs the
reductions and uses Jacobian doubling, mixed addition and left-to-right
binary multiplication. It checks every proper-divisor order certificate,
the Chinese remainder congruences, trial-division primality, and the
order-five reduction used in Lemma~\ref{lem:nontorsion}. Exact arithmetic
in $\Q(\sqrt2)$ authenticates $P=2R_0$. Laurent-polynomial arithmetic
checks the translation and sign identities.

Run
\begin{verbatim}
python -B negative_results/computation/mixed_sign_v42/build.py
python -B negative_results/computation/mixed_sign_v42/check.py
\end{verbatim}
The checker also accepts a certificate path via \texttt{--certificate}.
The finite computations establish the displayed group identities;
the progression argument above supplies the infinitely many
characteristic-zero examples.
\end{document}
